% ORIG03-SEGMENT-CODE: OR03-S005
% ORIG03-SEGMENT-ID: d90.orig.v1.tr03.seg0005
% ORIG03-MODULE-ID: ps03.module.0001
% SPDX-License-Identifier: CC-BY-SA-4.0
% Karya asli berbahasa Indonesia; semua butir disusun independen.

\section{Set soal III: dualitas dan kondisi KKT}
\label{orig03:section:set-soal-dualitas-kkt}

Lima masalah berikut melatih pembentukan fungsi dual, kondisi
Karush--Kuhn--Tucker (KKT), kondisi Slater, dan sensitivitas nilai optimal.

% ORIG03-STABLE-ID: ps03.problem.0001
\begin{exercise}[Kendala persamaan skalar]
\label{orig03:ps03:problem:0001}
Pertimbangkan masalah
\[
 \min_{x\in\mathbb{R}}\ \frac12x^2
 \qquad\text{dengan kendala}\qquad x=1.
\]
\begin{enumerate}
% ORIG03-STABLE-ID: ps03.prompt.0001
\item\label{orig03:ps03:prompt:0001}
Turunkan fungsi dual, lalu tentukan pengali dan nilai dual optimal.
% ORIG03-STABLE-ID: ps03.prompt.0002
\item\label{orig03:ps03:prompt:0002}
Verifikasi semua kondisi KKT dan kesenjangan dualitas nol.
\end{enumerate}
\end{exercise}

% ORIG03-STABLE-ID: ps03.hint1.0001
\par\noindent\textbf{Petunjuk tahap 1 (ps03.prompt.0001).}
Gunakan Lagrangian $L(x,\lambda)=\tfrac12x^2+\lambda(x-1)$.
\label{orig03:ps03:hint1:0001}
% ORIG03-STABLE-ID: ps03.hint2.0001
\par\noindent\textbf{Petunjuk tahap 2 (ps03.prompt.0001).}
Infimum terhadap $x$ tercapai di $x=-\lambda$.
\label{orig03:ps03:hint2:0001}
% ORIG03-STABLE-ID: ps03.answer.0001
\par\noindent\textbf{Jawaban ringkas (ps03.prompt.0001).}
$g(\lambda)=-\tfrac12\lambda^2-\lambda$,
$\lambda^\star=-1$, dan $d^\star=1/2$.
\label{orig03:ps03:answer:0001}
% ORIG03-STABLE-ID: ps03.solution.0001
\par\noindent\textbf{Solusi lengkap (ps03.prompt.0001).}
Meminimumkan Lagrangian terhadap $x$ memberi syarat $x+\lambda=0$, jadi
$x=-\lambda$. Karena Lagrangian konveks kuat terhadap $x$, titik ini memberi
infimum dan
\[
 g(\lambda)=L(-\lambda,\lambda)
 =-\frac12\lambda^2-\lambda.
\]
Masalah dual memaksimalkan kuadratik cekung ini pada
$\lambda\in\mathbb{R}$. Persamaan $g'(\lambda)=-\lambda-1=0$ memberi
$\lambda^\star=-1$ dan $g(-1)=1/2$.
\emph{Rubrik bukti: \ref{orig03:rubric:proof:0006}.}
\label{orig03:ps03:solution:0001}

% ORIG03-STABLE-ID: ps03.hint1.0002
\par\noindent\textbf{Petunjuk tahap 1 (ps03.prompt.0002).}
KKT untuk kendala persamaan terdiri atas kelayakan primal dan stasioneritas.
\label{orig03:ps03:hint1:0002}
% ORIG03-STABLE-ID: ps03.hint2.0002
\par\noindent\textbf{Petunjuk tahap 2 (ps03.prompt.0002).}
Substitusikan $x^\star=1$ dan $\lambda^\star=-1$.
\label{orig03:ps03:hint2:0002}
% ORIG03-STABLE-ID: ps03.answer.0002
\par\noindent\textbf{Jawaban ringkas (ps03.prompt.0002).}
$x^\star=1$ layak dan
$x^\star+\lambda^\star=0$; nilai primal dan dual sama-sama $1/2$.
\label{orig03:ps03:answer:0002}
% ORIG03-STABLE-ID: ps03.solution.0002
\par\noindent\textbf{Solusi lengkap (ps03.prompt.0002).}
Kelayakan primal memberi $x^\star-1=0$. Stasioneritas memberi
$\nabla_xL(x^\star,\lambda^\star)=x^\star+\lambda^\star=1-1=0$.
Tidak ada syarat tanda atau kelengkapan komplementer untuk pengali persamaan.
Nilai primal adalah $p^\star=\tfrac12(1)^2=1/2$, sedangkan hasil dual pada
subbutir sebelumnya memberi $d^\star=1/2$. Jadi
$p^\star-d^\star=0$.
\emph{Rubrik bukti: \ref{orig03:rubric:proof:0006}.}
\label{orig03:ps03:solution:0002}

% ORIG03-STABLE-ID: ps03.problem.0002
\begin{exercise}[Kendala ketaksamaan aktif]
\label{orig03:ps03:problem:0002}
Pertimbangkan
\[
 \min_{x\in\mathbb{R}}\ (x-2)^2
 \qquad\text{dengan kendala}\qquad x\leq1.
\]
\begin{enumerate}
% ORIG03-STABLE-ID: ps03.prompt.0003
\item\label{orig03:ps03:prompt:0003}
Gunakan kondisi KKT untuk menentukan solusi primal dan pengali optimal.
% ORIG03-STABLE-ID: ps03.prompt.0004
\item\label{orig03:ps03:prompt:0004}
Verifikasi kondisi Slater, turunkan fungsi dual, dan periksa nilainya di
pengali optimal.
\end{enumerate}
\end{exercise}

% ORIG03-STABLE-ID: ps03.hint1.0003
\par\noindent\textbf{Petunjuk tahap 1 (ps03.prompt.0003).}
Tuliskan kendala sebagai $x-1\leq0$ dengan pengali $\lambda\geq0$.
\label{orig03:ps03:hint1:0003}
% ORIG03-STABLE-ID: ps03.hint2.0003
\par\noindent\textbf{Petunjuk tahap 2 (ps03.prompt.0003).}
Gunakan stasioneritas $2(x-2)+\lambda=0$ dan
kelengkapan komplementer $\lambda(x-1)=0$.
\label{orig03:ps03:hint2:0003}
% ORIG03-STABLE-ID: ps03.answer.0003
\par\noindent\textbf{Jawaban ringkas (ps03.prompt.0003).}
$x^\star=1$ dan $\lambda^\star=2$.
\label{orig03:ps03:answer:0003}
% ORIG03-STABLE-ID: ps03.solution.0003
\par\noindent\textbf{Solusi lengkap (ps03.prompt.0003).}
Jika kendala tidak aktif, kelengkapan komplementer memberi $\lambda=0$ dan
stasioneritas memberi $x=2$, yang tidak layak. Maka kendala aktif:
$x^\star=1$. Stasioneritas kemudian memberi
$2(1-2)+\lambda^\star=0$, sehingga $\lambda^\star=2\geq0$.
Kelayakan primal, kelayakan dual, stasioneritas, dan kelengkapan komplementer
semuanya terpenuhi. Karena masalah konveks, titik KKT ini optimal.
\emph{Rubrik bukti: \ref{orig03:rubric:proof:0006}.}
\label{orig03:ps03:solution:0003}

% ORIG03-STABLE-ID: ps03.hint1.0004
\par\noindent\textbf{Petunjuk tahap 1 (ps03.prompt.0004).}
$x=0$ adalah kandidat titik Slater.
\label{orig03:ps03:hint1:0004}
% ORIG03-STABLE-ID: ps03.hint2.0004
\par\noindent\textbf{Petunjuk tahap 2 (ps03.prompt.0004).}
Dalam $L(x,\lambda)=(x-2)^2+\lambda(x-1)$, infimum terhadap $x$ tercapai
di $x=2-\lambda/2$.
\label{orig03:ps03:hint2:0004}
% ORIG03-STABLE-ID: ps03.answer.0004
\par\noindent\textbf{Jawaban ringkas (ps03.prompt.0004).}
Slater berlaku; $g(\lambda)=\lambda-\lambda^2/4$ untuk $\lambda\geq0$, dan
$g(2)=1=p^\star$.
\label{orig03:ps03:answer:0004}
% ORIG03-STABLE-ID: ps03.solution.0004
\par\noindent\textbf{Solusi lengkap (ps03.prompt.0004).}
Titik $x=0$ memenuhi $x-1=-1<0$, jadi kondisi Slater berlaku. Untuk
$\lambda\geq0$, peminimum Lagrangian terhadap $x$ adalah
$x(\lambda)=2-\lambda/2$. Substitusi memberi
\[
 g(\lambda)=\left(-\frac\lambda2\right)^2
 +\lambda\left(1-\frac\lambda2\right)
 =\lambda-\frac{\lambda^2}{4}.
\]
Kuadratik cekung ini maksimum di $\lambda=2$ dengan nilai $1$. Nilai primal
di $x^\star=1$ juga $(1-2)^2=1$, sesuai dualitas kuat yang dijamin Slater.
\emph{Rubrik bukti: \ref{orig03:rubric:proof:0006}.}
\label{orig03:ps03:solution:0004}

% ORIG03-STABLE-ID: ps03.problem.0003
\begin{exercise}[Proyeksi norma minimum pada hiperbidang]
\label{orig03:ps03:problem:0003}
Ambil $a\in\mathbb{R}^n\setminus\{0\}$ dan $b\in\mathbb{R}$, lalu
pertimbangkan
\[
 \min_x\ \frac12\lVert x\rVert_2^2
 \qquad\text{dengan kendala}\qquad a^\top x=b.
\]
\begin{enumerate}
% ORIG03-STABLE-ID: ps03.prompt.0005
\item\label{orig03:ps03:prompt:0005}
Selesaikan sistem KKT untuk memperoleh $x^\star$ dan $\lambda^\star$.
% ORIG03-STABLE-ID: ps03.prompt.0006
\item\label{orig03:ps03:prompt:0006}
Turunkan fungsi dual dan verifikasi bahwa nilai primal dan dual optimal sama.
% ORIG03-STABLE-ID: ps03.prompt.0007
\item\label{orig03:ps03:prompt:0007}
Buktikan bahwa solusi primal itu unik dan bahwa kondisi KKT cukup untuk
optimalitas pada masalah ini.
\end{enumerate}
\end{exercise}

% ORIG03-STABLE-ID: ps03.hint1.0005
\par\noindent\textbf{Petunjuk tahap 1 (ps03.prompt.0005).}
Gunakan $L(x,\lambda)=\tfrac12\lVert x\rVert_2^2+
\lambda(a^\top x-b)$.
\label{orig03:ps03:hint1:0005}
% ORIG03-STABLE-ID: ps03.hint2.0005
\par\noindent\textbf{Petunjuk tahap 2 (ps03.prompt.0005).}
Stasioneritas memberi $x=-\lambda a$; substitusikan ke kendala.
\label{orig03:ps03:hint2:0005}
% ORIG03-STABLE-ID: ps03.answer.0005
\par\noindent\textbf{Jawaban ringkas (ps03.prompt.0005).}
\[
 x^\star=\frac{b}{\lVert a\rVert_2^2}a,
 \qquad
 \lambda^\star=-\frac{b}{\lVert a\rVert_2^2}.
\]
\label{orig03:ps03:answer:0005}
% ORIG03-STABLE-ID: ps03.solution.0005
\par\noindent\textbf{Solusi lengkap (ps03.prompt.0005).}
Stasioneritas adalah $x+\lambda a=0$, sehingga $x=-\lambda a$.
Kelayakan persamaan lalu memberi
$-\lambda\lVert a\rVert_2^2=b$. Karena $a\neq0$, penyebutnya positif dan
\[
 \lambda^\star=-\frac{b}{\lVert a\rVert_2^2},
 \qquad
 x^\star=-\lambda^\star a=
 \frac{b}{\lVert a\rVert_2^2}a.
\]
\emph{Rubrik bukti: \ref{orig03:rubric:proof:0006}.}
\label{orig03:ps03:solution:0005}

% ORIG03-STABLE-ID: ps03.hint1.0006
\par\noindent\textbf{Petunjuk tahap 1 (ps03.prompt.0006).}
Ambil infimum Lagrangian di $x=-\lambda a$.
\label{orig03:ps03:hint1:0006}
% ORIG03-STABLE-ID: ps03.hint2.0006
\par\noindent\textbf{Petunjuk tahap 2 (ps03.prompt.0006).}
Maksimalkan kuadratik cekung yang diperoleh terhadap $\lambda\in\mathbb{R}$.
\label{orig03:ps03:hint2:0006}
% ORIG03-STABLE-ID: ps03.answer.0006
\par\noindent\textbf{Jawaban ringkas (ps03.prompt.0006).}
$g(\lambda)=-\tfrac12\lambda^2\lVert a\rVert_2^2-\lambda b$ dan
$p^\star=d^\star=b^2/(2\lVert a\rVert_2^2)$.
\label{orig03:ps03:answer:0006}
% ORIG03-STABLE-ID: ps03.solution.0006
\par\noindent\textbf{Solusi lengkap (ps03.prompt.0006).}
Substitusi $x=-\lambda a$ ke Lagrangian memberi
\[
 g(\lambda)=-\frac12\lambda^2\lVert a\rVert_2^2-\lambda b.
\]
Turunannya nol pada
$\lambda^\star=-b/\lVert a\rVert_2^2$, dan substitusi memberi
$d^\star=b^2/(2\lVert a\rVert_2^2)$. Di sisi primal,
\[
 \frac12\lVert x^\star\rVert_2^2
 =\frac12\frac{b^2\lVert a\rVert_2^2}{\lVert a\rVert_2^4}
 =\frac{b^2}{2\lVert a\rVert_2^2}=d^\star.
\]
\label{orig03:ps03:solution:0006}

% ORIG03-STABLE-ID: ps03.hint1.0007
\par\noindent\textbf{Petunjuk tahap 1 (ps03.prompt.0007).}
Gunakan konveksitas kuat objektif untuk keunikan.
\label{orig03:ps03:hint1:0007}
% ORIG03-STABLE-ID: ps03.hint2.0007
\par\noindent\textbf{Petunjuk tahap 2 (ps03.prompt.0007).}
Gunakan ketaksamaan konveksitas orde pertama pada titik KKT dan hilangkan
suku linear dengan kelayakan persamaan.
\label{orig03:ps03:hint2:0007}
% ORIG03-STABLE-ID: ps03.answer.0007
\par\noindent\textbf{Jawaban ringkas (ps03.prompt.0007).}
Objektif konveks kuat pada himpunan layak konveks, sehingga peminimum tunggal;
stasioneritas dan kelayakan KKT menjamin optimalitas global.
\label{orig03:ps03:answer:0007}
% ORIG03-STABLE-ID: ps03.solution.0007
\par\noindent\textbf{Solusi lengkap (ps03.prompt.0007).}
Fungsi $f(x)=\tfrac12\lVert x\rVert_2^2$ bersifat konveks kuat dengan
parameter satu. Pembatasannya pada hiperbidang afin tetap konveks kuat, jadi
tidak mungkin mempunyai dua peminimum berbeda. Untuk kecukupan, ambil
pasangan KKT $(x^\star,\lambda^\star)$. Stasioneritas memberi
$\nabla f(x^\star)=-\lambda^\star a$. Untuk setiap $x$ layak, konveksitas
memberi
\[
 f(x)\geq f(x^\star)+\nabla f(x^\star)^\top(x-x^\star)
 =f(x^\star)-\lambda^\star a^\top(x-x^\star)=f(x^\star),
\]
karena $a^\top x=a^\top x^\star=b$. Maka titik KKT optimal global, dan
konveksitas kuat membuatnya unik.
\emph{Rubrik bukti: \ref{orig03:rubric:proof:0006}.}
\label{orig03:ps03:solution:0007}

% ORIG03-STABLE-ID: ps03.problem.0004
\begin{exercise}[Sensitivitas ruas kanan]
\label{orig03:ps03:problem:0004}
Definisikan fungsi nilai
\[
 p(u)=\min_x\left\{\frac12x^2:x=u\right\}.
\]
\begin{enumerate}
% ORIG03-STABLE-ID: ps03.prompt.0008
\item\label{orig03:ps03:prompt:0008}
Hitung $p(u)$, pengali optimal $\lambda^\star(u)$ untuk kendala $x-u=0$,
dan hubungan antara $p'(u)$ dan pengali itu.
% ORIG03-STABLE-ID: ps03.prompt.0009
\item\label{orig03:ps03:prompt:0009}
Untuk perturbasi $\delta$, turunkan ekspansi eksak $p(u+\delta)$ dan
identifikasi suku sensitivitas orde pertama serta galat orde keduanya.
\end{enumerate}
\end{exercise}

% ORIG03-STABLE-ID: ps03.hint1.0008
\par\noindent\textbf{Petunjuk tahap 1 (ps03.prompt.0008).}
Kendala langsung memaksa $x=u$.
\label{orig03:ps03:hint1:0008}
% ORIG03-STABLE-ID: ps03.hint2.0008
\par\noindent\textbf{Petunjuk tahap 2 (ps03.prompt.0008).}
Stasioneritas Lagrangian
$\tfrac12x^2+\lambda(x-u)$ adalah $x+\lambda=0$.
\label{orig03:ps03:hint2:0008}
% ORIG03-STABLE-ID: ps03.answer.0008
\par\noindent\textbf{Jawaban ringkas (ps03.prompt.0008).}
$p(u)=u^2/2$, $\lambda^\star(u)=-u$, dan
$p'(u)=u=-\lambda^\star(u)$.
\label{orig03:ps03:answer:0008}
% ORIG03-STABLE-ID: ps03.solution.0008
\par\noindent\textbf{Solusi lengkap (ps03.prompt.0008).}
Satu-satunya titik layak adalah $x^\star(u)=u$, sehingga
$p(u)=u^2/2$. Stasioneritas memberi
$x^\star(u)+\lambda^\star(u)=0$, jadi $\lambda^\star(u)=-u$.
Mendiferensialkan fungsi nilai menghasilkan
$p'(u)=u=-\lambda^\star(u)$. Tanda minus muncul karena kendala ditulis
sebagai $x-u=0$.
\emph{Rubrik bukti: \ref{orig03:rubric:proof:0007}.}
\label{orig03:ps03:solution:0008}

% ORIG03-STABLE-ID: ps03.hint1.0009
\par\noindent\textbf{Petunjuk tahap 1 (ps03.prompt.0009).}
Kembangkan $(u+\delta)^2/2$.
\label{orig03:ps03:hint1:0009}
% ORIG03-STABLE-ID: ps03.hint2.0009
\par\noindent\textbf{Petunjuk tahap 2 (ps03.prompt.0009).}
Ganti $u$ pada suku linear dengan $-\lambda^\star(u)$.
\label{orig03:ps03:hint2:0009}
% ORIG03-STABLE-ID: ps03.answer.0009
\par\noindent\textbf{Jawaban ringkas (ps03.prompt.0009).}
$p(u+\delta)=p(u)-\lambda^\star(u)\delta+\delta^2/2$.
\label{orig03:ps03:answer:0009}
% ORIG03-STABLE-ID: ps03.solution.0009
\par\noindent\textbf{Solusi lengkap (ps03.prompt.0009).}
Perhitungan eksak memberi
\[
 p(u+\delta)=\frac12(u+\delta)^2
 =\frac12u^2+u\delta+\frac12\delta^2
 =p(u)-\lambda^\star(u)\delta+\frac12\delta^2.
\]
Jadi $-\lambda^\star(u)\delta$ adalah perubahan linear yang diprediksi
pengali, sedangkan $\delta^2/2$ adalah galat orde kedua yang eksak.
\emph{Rubrik bukti: \ref{orig03:rubric:proof:0007}.}
\label{orig03:ps03:solution:0009}

% ORIG03-STABLE-ID: ps03.problem.0005
\begin{exercise}[Proyeksi pada bola Euclid melalui KKT]
\label{orig03:ps03:problem:0005}
Ambil $c\in\mathbb{R}^n$ dan $r$ dengan $0<r<\lVert c\rVert_2$.
Pertimbangkan
\[
 \min_x\ \frac12\lVert x-c\rVert_2^2
 \qquad\text{dengan kendala}\qquad \lVert x\rVert_2\leq r.
\]
Gunakan representasi mulus
$g(x)=\tfrac12(\lVert x\rVert_2^2-r^2)\leq0$.
\begin{enumerate}
% ORIG03-STABLE-ID: ps03.prompt.0010
\item\label{orig03:ps03:prompt:0010}
Selesaikan kondisi KKT untuk menentukan $x^\star$ dan pengali
$\mu^\star$.
% ORIG03-STABLE-ID: ps03.prompt.0011
\item\label{orig03:ps03:prompt:0011}
Verifikasi kondisi Slater dan buktikan bahwa pasangan KKT tersebut merupakan
solusi primal-dual global yang unik pada sisi primal.
\end{enumerate}
\end{exercise}

% ORIG03-STABLE-ID: ps03.hint1.0010
\par\noindent\textbf{Petunjuk tahap 1 (ps03.prompt.0010).}
Stasioneritas adalah $x-c+\mu x=0$.
\label{orig03:ps03:hint1:0010}
% ORIG03-STABLE-ID: ps03.hint2.0010
\par\noindent\textbf{Petunjuk tahap 2 (ps03.prompt.0010).}
Karena $c$ berada di luar bola, kendala optimal aktif; gunakan
$x=c/(1+\mu)$ dan $\lVert x\rVert_2=r$.
\label{orig03:ps03:hint2:0010}
% ORIG03-STABLE-ID: ps03.answer.0010
\par\noindent\textbf{Jawaban ringkas (ps03.prompt.0010).}
\[
 x^\star=r\frac{c}{\lVert c\rVert_2},
 \qquad
 \mu^\star=\frac{\lVert c\rVert_2}{r}-1>0.
\]
\label{orig03:ps03:answer:0010}
% ORIG03-STABLE-ID: ps03.solution.0010
\par\noindent\textbf{Solusi lengkap (ps03.prompt.0010).}
Lagrangiannya
$L(x,\mu)=\tfrac12\lVert x-c\rVert_2^2+
\tfrac\mu2(\lVert x\rVert_2^2-r^2)$ dengan $\mu\geq0$.
Stasioneritas memberi $(1+\mu)x=c$. Jika $\mu=0$, maka $x=c$, yang tidak
layak karena $\lVert c\rVert_2>r$. Jadi $\mu^\star>0$, dan kelengkapan
komplementer memaksa $\lVert x^\star\rVert_2=r$. Dari
$x^\star=c/(1+\mu^\star)$ diperoleh
$1+\mu^\star=\lVert c\rVert_2/r$, sehingga rumus yang dinyatakan berlaku.
\emph{Rubrik bukti: \ref{orig03:rubric:proof:0006}.}
\label{orig03:ps03:solution:0010}

% ORIG03-STABLE-ID: ps03.hint1.0011
\par\noindent\textbf{Petunjuk tahap 1 (ps03.prompt.0011).}
$x=0$ memenuhi kendala secara ketat.
\label{orig03:ps03:hint1:0011}
% ORIG03-STABLE-ID: ps03.hint2.0011
\par\noindent\textbf{Petunjuk tahap 2 (ps03.prompt.0011).}
Gabungkan Slater, konveksitas masalah, dan konveksitas kuat objektif.
\label{orig03:ps03:hint2:0011}
% ORIG03-STABLE-ID: ps03.answer.0011
\par\noindent\textbf{Jawaban ringkas (ps03.prompt.0011).}
Slater berlaku di $x=0$; KKT perlu dan cukup, dualitas kuat berlaku, dan
$x^\star$ unik karena objektif konveks kuat.
\label{orig03:ps03:answer:0011}
% ORIG03-STABLE-ID: ps03.solution.0011
\par\noindent\textbf{Solusi lengkap (ps03.prompt.0011).}
Karena $r>0$, titik $x=0$ memenuhi
$g(0)=-r^2/2<0$, jadi kondisi Slater berlaku. Objektif dan fungsi kendala
konveks serta terdiferensialkan; karena itu, di bawah Slater, kondisi KKT
bersifat perlu dan cukup dan tidak ada kesenjangan dualitas. Pasangan pada
subbutir sebelumnya memenuhi kelayakan primal, $\mu^\star\geq0$,
stasioneritas, dan kelengkapan komplementer, sehingga optimal global.
Objektif $\tfrac12\lVert x-c\rVert_2^2$ bersifat konveks kuat, maka
peminimum primal pada himpunan bola yang konveks adalah unik.
\emph{Rubrik bukti: \ref{orig03:rubric:proof:0006}.}
\label{orig03:ps03:solution:0011}
